\section{Introducción: la socialización con IA no es «una IA que platica contigo»}

Esta entrevista gira en torno a dos productos de Ziran Xuanze (自然选择): el juego con IA Eve y la red social con IA Elys. La narrativa central que plantea Tristan es que «los seres humanos están demasiado solos», pero no se trata de un simple producto de acompañamiento emocional. La ambición mayor de Elys es que el doble de IA salga activamente al mundo y ayude al usuario a convertir su propio context en conexiones reales. Intenta responder una pregunta que dejó pendiente la era de internet: por qué las redes sociales siempre comprimen a personas complejas en etiquetas de baja dimensión, y si la IA puede, por primera vez, hacer que fluya el context de alta dimensión.

\begin{importantbox}{Tesis central de este episodio}
La socialización tradicional en internet depende de etiquetas de baja dimensión y de las acciones que el propio usuario emprende; las redes sociales con IA intentan que el usuario construya su subjetividad, que su doble cibernético, con su context a cuestas, se exprese, recomiende y rompa el hielo por iniciativa propia, y que al final conduzca de vuelta a conexiones entre personas reales.
\end{importantbox}

\begin{knowledgebox}{Nota sobre la estrategia visual}
Este video es una entrevista con encuadre fijo, sin slides didácticas, pizarrón ni demostraciones del producto. Conforme al estándar de este repositorio para pódcasts, el cuerpo del texto no repite fotogramas de las personas; la portada sirve para identificar la fuente, y el cuerpo del texto presenta la lógica del producto mediante diagramas conceptuales y tablas.
\end{knowledgebox}

\subsection{Resumen de la sección}

Elys no es un simple Tinder con IA, ni tampoco un juguete en el que una IA platica con otra IA. Sus variables clave son el context, la subjetividad, el intercambio de privacidad, la tasa de conexión entre personas reales y la forma de producto proactive.
